\section{Related Work}

Previous studies on tensor networks in machine learning have mainly been devoted to analyzing the theoretical properties of neural networks. A better understanding of feed-forward, convolutional and recurrent architectures has been gained, including compression parameters \citep{novikov2015tensorizing}, expressive power \citep{cohen2016expressive, cohen2016convolutional, khrulkov2018expressive}, and depth efficiency for long-term memory \citep{levine2018benefits}. For sequence modeling tasks in NLP, there are two stages of the previous research.

\paragraph{Theoretical Proposals.}
\citep{pestun2017tensor} propose a tensor network language model aims to construct the long-range correlation in real-world language modeling, \citep{pestun2017language} propose a quantum statistical language model on a one-dimensional lattice which is called trace-density model. 
To the best of our knowledge, these tensor network language models have remained a theoretical proposal instead of an empirical one. 


\paragraph{Sequence Modeling.} \cite{novikov2021tensor} propose a new efficient tensor train-based approach to tensor-train density estimation that allows efficient computation of probability density function. \cite{miller2021tensor} apply a recurrent tensor network, uniform matrix product state, to the probabilistic sequence modeling while opening significant new research directions in the design of sequential generative models. However, probably due to the efficiency issue, most of the existing models have only been tested on small vocabulary size and simple datasets like Tomita grammars \citep{tomita1982dynamic}, further evaluation on moderately-scaled natural language datasets is necessary to fully assess their performance. We provide a comparison of the datasets we used with related work in the Appendix: \ref{sec:datasets}.

This paper is the first work to  derive a tensor network language model in a way that can be applied to real-world language modeling datasets. The efforts to improve efficiency are twofold. First, we do not calculate the probability normalization term used for the total probability law; instead, we turn to  calculate conditional probabilities based on context representations, as described in Sec. \ref{sec:recursive probability}.
Second, we further decompose TT cores and use  low-scale hidden units, see in Sec. \ref{sec: Tinyttlm}.


\section{预备知识}
\label{sec:preliminaries}
我们在此简要回顾一些基本概念与记号\footnote{此处的大部分记号沿用教材 \textit{Deep Learning}
\cite{goodfellow2016deep}。}；完整的技术性介绍可参阅标准教科书 \cite{BI20221, 10.5555/1643460}。

% \textbf{Tensor} is a multidimensional array representing a multilinear relationship between certain objects in vectors spaces. First order tensors vectors denoted $\mathbf{v}$. Second order tensors are matrices, denoted $\mathbf{M}$. By $\mathcal{T}$, we denote tensors of order 3 or greater. For a third order $\mathcal{T}$, we denote its element $(i,j,k)$ as $\mathcal{T}_{ijk}$.

\paragraph{记号。}在本文中，每个\textit{张量（tensor）} $\tA$ 都是元素（称为\emph{分量（component）}）取自 $\mathbb{R}$ 的多维数组（multi-dimensional array），每个元素由其在数组中的整数坐标标记；例如，对于二维数组，位于 $i,j \in \mathbb{N}$ 处的分量记作 $\etA_{ij}$。
 张量的\emph{阶（order）}指它具有多少个指标（例如，向量 $\vv$ 是一阶张量，矩阵 $\mM$ 是二阶张量，等等）。张量的\emph{维数（dimension）}指某个特定指标（即所谓的\textit{模态（mode）}）可取值的数目，例如，$\tB \in \mathbb{R}^{I_1 \times I_2 \times I_3}$ 的维数为 $I_1 \times I_2 \times I_3$。

\paragraph{张量积 \cite{cohen2016expressive}。}对于两个张量 $\tC \in \mathbb{R}^{I_1 \times \cdots \times I_j}$（$j$ 阶）与 $\tD \in \mathbb{R}^{I_{j+1}, \times \cdots \times I_{j+k}}$（$k$ 阶），它们的\textit{张量积（tensor product）}记作 $\otimes$，返回一个张量 $\etE_{i_1 \cdots i_{j+k}} = \etC_{i_1...i_j} \cdot \etD_{i_{j+1} \cdots i_{j+k}}$（$j+k$ 阶）。注意，当 $j = k = 1$ 时，张量积退化为向量之间的外积（outer product）。

\paragraph{广义内积 \citep{kossaifi2020tensor}。}
\label{Generalized inner product}
对于尺寸相同的两个张量 $\tX,\tY\in\mathbb{R}^{I_1\times I_2 \times \cdots \times I_N}$，它们的\textit{内积（inner product）}定义为 $\langle \tX,\tY \rangle=\sum_{i_1=1}^{I_1}\sum_{i_2=1}^{I_2} \cdots \sum_{i_N=1}^{I_N} \etX_{i_1,i_2,...,i_N} \etY_{i_1,i_2,...,i_N}$。对于共享 $N$ 个同尺寸模态（mode）的两个张量 $\tX \in \mathbb{R}^{I_1\times I_2 \times \cdots \times I_N \times I_x}$ 与 $\tY \in \mathbb{R}^{I_1\times I_2\times \cdots I_N \times I_y}$，“广义内积（generalized inner product）”计算为
%
\begin{align}
    \langle \tX,\tY \rangle_N =\sum_{i_1=1}^{I_1}\sum_{i_2=1}^{I_2}\cdot\cdot\cdot\sum_{i_N=1}^{I_N} \etX_{i_1,i_2,...,i_N} \etY_{i_1,i_2,...,i_N} \nonumber
\end{align}
%
且 $\langle \tX,\tY \rangle_N \in \mathbb{R}^{I_x \times I_y}$。



\begin{figure*}[t]
\centering
% \vspace{-15pt}
\includegraphics[scale=0.45]{figure/preliminary.pdf}
\caption{\textit{张量图示记法（tensor diagram notation）}简介。张量图有两条规则：(1) 张量用实心图形表示，图形所带“腿”的数目对应其指标的数目；(2) 连接两条指标线意味着对所连指标进行\textit{收缩（contraction）}或求和。在本文中，我们在公式旁附上这些图示，以便于理解。 }
\label{fig: tensor_diagram}
\end{figure*}


